\subsection{At the Threshold}\label{sec:threshold}

On 24 December the Church opens the Mass of the Vigil with words from Exodus,
\latin{Hodie scietis} (gloss: this day you shall know), and with them the waiting of four weeks
reaches its last day. The O antiphons have run their course, the Ember days are past, and by None
the season is over.

The Collect of the Vigil (Section~\ref{sec:last-days}) is the prayer for that day. On the Vigil
itself, the rule of life can give way to preparation of another kind: the house made
ready, the confession made if it has not been, the household gathered. If there is time in the
afternoon, read the Gospel of the Vigil, the angel's message to Joseph (Matthew 1:18--21), and say
the Collect once more. Say the Angelus at evening with whoever is at home.

St Bernard told his monks to ponder how much they had received by the first coming and how much
they were promised at the second (Section~\ref{sec:three-comings}). The night of the Nativity
remembers the first. The Christian who keeps it goes on afterward as Advent has taught: receiving
him in the word and the sacraments, and looking for his coming in glory.
